DPitG resolves the tension inherent in PitG --- requiring both precision and
decisiveness simultaneously --- at a modest median sampling cost
(Sections~\ref{sec:fair_coin_large_scale}--\ref{sec:trends}).
We now place this result in methodological context, draw out its practical
implications, and candidly acknowledge its limitations.

DPitG operates under a different philosophy from classical sequential methods.
Group sequential designs \citep{jennison2000} control the familywise Type~I error rate
across interim analyses via alpha-spending functions, and are therefore tied to the
number and timing of planned looks.
DPitG requires no such correction, since the Bayesian posterior accounts for all
accumulated data without reference to the stopping intention.
This property extends to simultaneous tests: the DPitG stopping rule applies
unchanged across independent parallel experiments without multiple-comparisons
correction; Appendix~\ref{app:multiple_comparisons} develops this in full,
together with all three types of Bayesian power analysis (prospective,
retrospective, and replication).

Closest in spirit among Bayesian sequential designs is Bayes Factor Design Analysis
(BFDA; \citealt{schonbrodt2018}), which plans sample size to achieve a target Bayes
Factor with specified probability.\footnote{This is a prospective planning use of the BF, distinct from using the BF as a
sequential stopping rule (see Appendix~\ref{app:bayes_factors} for a discussion
of BF stopping and its limitations).}
% \citet{schonbrodt2018}: Sch\"onbrodt \& Wagenmakers (2018) introduce BFDA as the
% natural Bayesian sequential counterpart to power analysis: it targets evidential
% strength (Bayes Factor) rather than posterior precision, making it complementary to
% rather than competing with DPitG.
The key distinction is that BFDA targets evidential strength (ratio of marginal
likelihoods) whereas DPitG targets posterior precision; the two are complementary
as planning tools, and a design requiring both a conclusive HDI+ROPE framework verdict and a
minimum Bayes Factor is a natural extension.
% [AUTHOR NOTE: standalone multiple-comparisons paragraph condensed to a pointer
%  after the no-correction sentence above; full content preserved below for reference.]
% The HDI+ROPE framework for the decision step offers a structural
% methodological advantage over frequentist approaches: unlike NHST
% (see Appendix~\ref{app:nhst}), where conducting
% multiple simultaneous tests inflates the Type~I error rate and requires corrections
% (e.g.\ Bonferroni) that reduce per-test power, the Bayesian posterior is unaffected ---
% $\omega_{\rm goal}$ and the DPitG stopping rule apply unchanged across independent
% simultaneous tests without correction \citep{gelman2012, kruschke2015doing}.
% \citet{gelman2012}: Gelman, Hill \& Yajima (2012) is the standard journal reference
% for the claim that Bayesian posteriors typically require no multiple-comparison
% correction when tests are based on independent models.
% \citet{kruschke2015doing}: DBDA2 (Chapter~11) provides the pedagogical treatment.
% For prospective sample-size planning, $N_{\rm goal}$
% (Equation~\ref{eq:pitg_stop_iteration}) provides a closed-form lower bound on the
% samples required to achieve $\omega_{\rm goal}$; DPitG may need additional samples
% when $\theta_{\rm true}$ lies near a ROPE boundary (see Section~\ref{sec:trends}).
% Appendix~\ref{app:multiple_comparisons} develops the multiple-comparisons point in
% full and covers all three types of Bayesian power analysis (prospective, retrospective,
% and replication).

While comparisons with competing sequential methods clarify DPitG's positioning,
its computational scope deserves equal attention.
DPitG is a generic sequential stopping rule applicable to any posterior admitting a
credible interval, including non-conjugate and hierarchical settings that require MCMC;
the present analysis restricts to conjugate posteriors, where the HDI is fully
analytical and no specialist probabilistic programming software is required.
% Tractable-posterior scope, case by case:
% - Single-group Bernoulli (this paper): conjugate Beta posterior, fully analytical
%   HDI; no sampling of any kind needed.
% - Single-group continuous (conjugate normal prior): normal or Student-t posterior,
%   fully analytical; no MCMC needed.
% - Two-group binomial: no closed form for the difference of two Betas, but
%   direct i.i.d. Monte Carlo suffices (draw from Beta_1, draw from Beta_2, subtract);
%   this is plain Monte Carlo, not a Markov chain.
% - Two-group continuous (weakly informative prior): by the Bernstein--von Mises
%   theorem the posterior converges to a normal, so Welch's standard error gives the
%   HDI width without sampling.
% - Non-conjugate informative priors: MCMC would be required; this case is deferred
%   to future work.
% - Hierarchical models: MCMC required; also out of scope here.


A recurring concern about PitG --- and by extension DPitG --- is that the precision
requirement can seem too demanding in practice.
Researchers frequently underestimate the noise in their measurements and are consequently surprised by the sample sizes that
$\omega_{\rm goal}$ implies (J.K.\ Kruschke, pers.\ comm.\ 2022). This is not a flaw
in the method but a form of epistemic honesty: imprecise posteriors produce inconclusive
results regardless of the algorithm applied; DPitG simply refuses to declare a verdict
until the posterior is genuinely informative. Setting $\omega_{\rm goal}$ thus forces
realistic planning rather than masking insufficient power behind a false reject (i.e., false positive) or a
silently inconclusive stop.
This discipline is a feature, not a limitation, for any analyst whose conclusions
must rest on genuinely informative data.

DPitG serves researchers in two complementary roles.
As an operational stopping rule, it is most naturally suited to settings where data
collection is relatively cheap and the budget comfortably exceeds $N_{\rm goal}$:
large-scale digital A/B testing \citep{johari2022}, opinion polling, and automated
quality-control pipelines are typical examples.
% \citet{johari2022}: Johari et al.\ (2022) study continuous monitoring of A/B tests ---
% a canonical high-volume, low-cost data-collection context where precision-based
% stopping is operationally natural and the budget constraint is rarely binding.
As a planning tool, it provides a diagnostic function even in resource-constrained
settings --- longitudinal studies, clinical trials, or expensive policy evaluations:
$N_{\rm goal}$ makes explicit what precision is achievable within the available budget,
enabling an honest assessment of what conclusions the data can and cannot support,
both prospectively (before collection begins) and retrospectively (to audit the
precision of a completed study).

The ROPE must be grounded in domain knowledge rather than chosen to suit the data
\citep{kruschke2018}; Appendix~\ref{app:practical_guidance} discusses principled
approaches for setting it --- including the MCID, bioequivalence standards, and
sensitivity analyses --- together with three broader considerations supporting
DPitG's adoption (equivalence testing, pre-registration, and the replication crisis).
Together with $\omega_{\rm goal}$ and $N_{\rm max}$, the ROPE completes the three
design parameters that must be fixed before data collection begins.

With those design inputs committed, the joint operating characteristics of DPitG can
be assessed prospectively: a researcher can check, analytically or via simulation,
that the planned $N_{\rm max}$ achieves the target conclusiveness rate at the effect
sizes of interest.
As the Results demonstrate, error rates are negligible except near the ROPE boundaries,
where neither an accept nor a reject verdict is easily justified.
A key practical implication (Section~\ref{sec:goal_impact}) is that conclusiveness
is more sensitive to $N_{\rm max}$ than to $\omega_{\rm goal}$: for the fair coin
case ($\Delta_{\rm ROPE}=0.1$), the median sampling premium relative to $N_{\rm goal}$
is zero for $\omega_{\rm goal}\le 0.07$ and rises to 41\% at $\omega_{\rm goal}=0.10$,
so setting $N_{\rm max}$ generously is the primary practical lever for ensuring high
conclusiveness across a range of precision goals.

We now turn to the limits of the present work.
Four limitations deserve acknowledgement.
First, the results are sensitive to the three design parameters --- the ROPE,
$\omega_{\rm goal}$, and $N_{\rm max}$ --- all of which must be committed to before
data collection begins.
A narrower ROPE (and with it a tighter $\omega_{\rm goal}$) increases the sample size
required for decisiveness and raises the
inconclusive rate near its boundaries.
Second, the stopping properties are validated empirically only for binary (Bernoulli)
data under a conjugate Beta prior.
The normal approximation underpinning $N_{\rm goal}$
(Equation~\ref{eq:pitg_stop_iteration}) may underestimate the required sample size
in small-sample settings or for highly skewed likelihoods; in such cases the exact
Binomial result in Appendix~\ref{app:expected_stop} should be preferred.
Third, the present analysis uses a flat (uniform) prior throughout, as a conservative
choice that avoids imposing unverified assumptions.
The flat prior contributes no pseudo-observations, requiring more data to achieve
$\omega_{\rm goal}$ than a well-justified informative prior would \citep[][Ch.~2]{gelman2013bayesian}.
% \citet{gelman2013bayesian}: BDA3 (Chapter~2) establishes the prior-as-pseudo-data
% interpretation for conjugate models: a Beta(a,b) prior is equivalent to having
% previously observed a successes and b failures, directly supporting the claim that
% informative priors contribute effective pseudo-observations.
Informative priors, where available and well-justified, can therefore serve as a
practical mitigation by narrowing the starting HDI width and lowering the required
sample size.
Conversely, a materially incorrect informative prior may bias the posterior in a way
that the precision criterion alone cannot detect \citep[][Ch.~6]{gelman2013bayesian};
% \citet{gelman2013bayesian}: BDA3 (Chapter~6) covers posterior predictive checks
% and prior sensitivity analysis --- the standard treatment for detecting and
% diagnosing prior misspecification.
we defer a full treatment of principled prior selection to a later study.
Fourth, while DPitG applies in principle to any posterior admitting such an interval,
this paper does not characterise its behaviour when the HDI is estimated from MCMC
samples.
In such settings the HDI estimate carries Monte Carlo error \citep{kruschke2015doing}
that could affect stopping
calibration --- particularly in early data collection when the posterior is diffuse ---
and this remains an open question for subsequent investigation.

% The extensions to continuous outcomes and two-group comparisons are
% theoretical derivations rather than empirically validated results.
% Confirming that the stopping properties demonstrated here for the Bernoulli case carry
% over --- in particular the reduction in inconclusive rates and the calibration of
% $N_{\rm stop}$ relative to $N_{\rm goal}$ --- requires a dedicated simulation study,
% which we flag as the primary empirical direction for future work.

